% =====================================================================
%  Apéndices A--G.  ESTADO: andamiaje.
%  Con `\appendix`, cada \chapter se numera como «Apéndice A, B, ...».
%  El contenido fuente está en memoria/apendices/*.md (legado a migrar).
% =====================================================================
\appendix
\label{cap:apendices}

\chapter{Manual técnico}
\pendienteredaccion{} Fuente: \texttt{memoria/apendices/A-manual-tecnico.md}.

\chapter{Paquete de replicación}
\pendienteredaccion{} Fuente: \texttt{memoria/apendices/B-replicacion.md}.

\chapter{Prompt RQ3 v1.0 y oráculo}
\pendienteredaccion{} Texto literal del prompt (\texttt{LlmPromptBuilder})
y descripción del oráculo. Fuente:
\texttt{memoria/apendices/C-prompt-rq3.md}.

\chapter{Catálogo de \texttt{DiscardReason}}
\pendienteredaccion{} Diez categorías (verificadas en
\texttt{DiscardReason.java}): \texttt{OUTER\_HAS\_ELSE},
\texttt{OUTER\_BLOCK\_MULTIPLE\_STATEMENTS},
\texttt{SINGLE\_STATEMENT\_NOT\_IF}, \texttt{INNER\_HAS\_ELSE},
\texttt{METHOD\_CALL\_IN\_CONDITION}, \texttt{ASSIGNMENT\_IN\_CONDITION},
\texttt{INCREMENT\_OR\_DECREMENT\_IN\_CONDITION},
\texttt{LAMBDA\_IN\_CONDITION}, \texttt{OBJECT\_CREATION\_IN\_CONDITION},
\texttt{NO\_NESTED\_IF\_PATTERN}. Fuente:
\texttt{memoria/apendices/D-discard-reasons.md}.

\chapter{Tablas extendidas RQ2 / RQ3}
\pendienteredaccion{} Fuente: \texttt{memoria/apendices/E-tablas-extendidas.md}.

\chapter{Procedimiento SonarQube}
\pendienteredaccion{} Fuente: \texttt{memoria/apendices/F-sonarqube.md} y
\texttt{docs/}.

\chapter{Trazabilidad RQ $\leftrightarrow$ artefactos}
\pendienteredaccion{} Fuente: \texttt{memoria/apendices/G-trazabilidad.md}.
